\section{Conclusões e Trabalho Futuro} \label{sec:concusoes-e-trabalho-futuro}

O presente trabalho alcançou o objetivo principal de desenvolver um sistema de \textit{Business Intelligence} focado na análise de incêndios florestais em Portugal Continental.
A dispersão de dados, identificada inicialmente como um desafio central, foi mitigada através da implementação de um processo ETL robusto, capaz de centralizar informação de natureza heterogénea.
A integração de dados operacionais do ICNF com fontes externas, nomeadamente a meteorologia (\textit{Open-Meteo}) e índices de vegetação, permitiu enriquecer o contexto analítico, oferecendo novas perspetivas sobre as ocorrências.

A adoção da metodologia de modelação dimensional de Kimball provou ser eficaz, facilitando a estruturação do \textit{Data Warehouse} e garantindo um desempenho adequado nas consultas. A definição clara dos processos de negócio (Ocorrência, Meteorologia, cobertura mediática, e Vegetação) permitiu criar um modelo intuitivo e alinhado com as necessidades de análise.

A etapa de \textit{Data Profiling} revelou-se crítica para a fiabilidade do sistema. A identificação de lacunas nas variáveis meteorológicas originais validou a decisão estratégica de recorrer a APIs externas, assegurando a completude da informação climática. Adicionalmente, a componente geográfica permitiu identificar padrões de concentração espacial (\textit{hotspots}), validando a utilidade do sistema para o planeamento estratégico e prevenção.

Em suma, a solução desenvolvida transforma dados brutos e dispersos em conhecimento acionável, oferecendo às entidades competentes uma ferramenta de apoio à decisão baseada em evidências históricas.

\subsection{Limitações do Estudo}
Apesar dos resultados positivos, o projeto apresenta limitações que devem ser consideradas.
A primeira reside na componente de cobertura mediática. Embora os dados tenham sido extraídos, a sua integração plena no modelo de análise carece ainda de técnicas avançadas de processamento de linguagem natural (\textit{NLP}) para aferir o sentimento e a relevância das notícias.

Uma segunda limitação prende-se com a escassez de informação recolhida sobre a vegetação antes e depois de cada incêndio, o que condiciona uma avaliação mais profunda do impacto ecológico e da recuperação dos solos.

Por fim, o período temporal analisado (2020-2025), embora recente e relevante, é relativamente curto para a identificação de tendências climáticas de longo prazo.

\subsection{Trabalho Futuro}
Para dar continuidade a este projeto e potenciar a sua utilidade, sugerem-se as seguintes linhas de investigação e desenvolvimento:

\begin{itemize}
	\item \textbf{Automatização integral do ETL:} Desenvolver mecanismos de extração automática para as fontes de Cobertura Mediática e Vegetação, eliminando a intervenção manual atualmente necessária e alinhando estes processos com a automação já existente para os dados do ICNF e Meteorologia;
	\item \textbf{Integração de Dados Operacionais (DECIR):} Implementar rotinas de extração de texto (\textit{OCR/Text Mining}) para processar os relatórios da Diretiva Operacional Nacional. Isto permitiria incorporar dados sobre meios aéreos e terrestres, possibilitando a criação de métricas sobre a eficácia do combate e tempos de resposta;
	\item \textbf{Análise de Sentimento (NLP):} Implementar algoritmos de processamento de linguagem natural para classificar automaticamente as notícias (positiva/negativa/neutra), permitindo correlacionar o impacto mediático com a severidade dos incêndios;
	\item \textbf{Integração em Tempo Real:} Adaptar a arquitetura para processamento em \textit{near real-time}, permitindo o acompanhamento de ocorrências ativas nos \textit{dashboards} de controlo;
	\item \textbf{Enriquecimento de Dados:} Incorporar dados sobre a orografia do terreno (declive e exposição solar), fatores determinantes na modelação da velocidade e direção de propagação das chamas.
\end{itemize}